% ============================================
% Johns Hopkins Letter Template
% Assistant Professor Cover Letter – Stony Brook University
% ============================================

\documentclass[11pt,letterpaper]{letter}

\usepackage{graphicx}
\usepackage{eso-pic}
\usepackage{geometry}
\usepackage{microtype}
\usepackage[hidelinks]{hyperref}

% ----- Page Setup -----
\geometry{left=25mm,right=25mm,top=25mm,bottom=25mm}

% ----- Logo Placement (first page only) -----
\newcommand{\PutJHULogoFG}[1]{%
  \AddToShipoutPictureFG*{%
    \AtPageUpperLeft{%
      \hspace{10mm}%
      \raisebox{-38mm}{%
        \includegraphics[width=6cm]{#1}%
      }%
    }%
  }%
}

% ----- Sender Info -----
\signature{Decory Edwards}
\address{Department of Economics \\ Johns Hopkins University \\ Baltimore, MD 21218 \\ \texttt{dedwar65@jh.edu}}

\begin{document}
\PutJHULogoFG{Images/jhulogo.png}

\begin{letter}{Search Committee \\
Department of Economics \\
Stony Brook University \\
Stony Brook, NY 11794 \\ USA}

\opening{Dear Members of the Search Committee,}

I am writing to apply for the position of Assistant Professor in the Department of Economics at Stony Brook University, beginning in Fall 2026. I am a Ph.D. candidate in Economics at Johns Hopkins University, advised by Professors Christopher D.~Carroll, Nicholas W.~Papageorge, and Stelios Fourakis, and I expect to receive my degree in May 2026. My research focuses on macroeconomics, inequality, and heterogeneous-agent modeling, with an emphasis on empirically grounded mechanisms that drive wealth and income dispersion.

My job-market paper examines heterogeneity in returns to wealth using micro-data from the Survey of Consumer Finances. I estimate the distribution and persistence of household-level portfolio returns and embed these estimates in a structural heterogeneous-agent framework. I show that allowing for persistent return heterogeneity substantially improves the model’s ability to match observed wealth concentration, offering a tractable way to connect micro-level financial frictions to aggregate inequality. A second project uses the Health and Retirement Study to study how interpersonal trust relates to realized portfolio returns, documenting a hump-shaped relationship that provides a behavioral channel linking attitudes to wealth accumulation.

Stony Brook’s commitment to rigorous empirical work, interdisciplinary collaboration, and research that speaks to national and global policy challenges closely matches the intellectual direction of my work. My focus on inequality, household heterogeneity, and macro-distributional dynamics complements the department’s strengths in applied microeconomics, public economics, behavioral economics, and computational methods. I am especially enthusiastic about contributing to a department that values research breadth, interdisciplinary engagement, and methodological innovation.

\textbf{Teaching.} At Johns Hopkins, I have taught and assisted in undergraduate macroeconomics and an upper-level macroeconomic policy seminar, leading weekly discussion sections and holding regular office hours. I emphasize clarity, reproducibility, and helping students build confidence through structured problem-solving and evidence-based reasoning. At Stony Brook, I would be prepared to teach undergraduate and graduate courses in macroeconomics, microeconomics, econometrics, and related field courses aligned with my research on inequality and household heterogeneity. I am committed to fostering inclusive, supportive learning environments that help students from diverse backgrounds succeed.

I believe I would be a strong fit for Stony Brook University because my research agenda (i) aligns with the department’s priorities in applied micro and inequality-related research, (ii) is well suited for interdisciplinary collaboration across the social sciences, and (iii) contributes to the university’s mission of advancing access, excellence, and innovation in public higher education. I am eager to contribute to your Ph.D. program, collaborate with faculty across research areas, and support the department’s teaching and service commitments.

Thank you for considering my application. I would be honored to join Stony Brook University’s Department of Economics and contribute to its research, teaching, and mentorship missions. I look forward to the opportunity to discuss my fit with the department.

\closing{Sincerely,}

\end{letter}
\end{document}
